\begin{lemma}[Claim 9: first tail use after the seed]
Fix the Section 7 minimum context and a
\noderef{HierarchyReplacementPresentation} \(P\).  Let \(\alpha\) be
missing on \(P\)'s last starred stage \(A\) and at \(y_\ell\).  Assume
the ambient Situation 7.1 condition
\(\noderef{HierarchyPathNumber}
(P.\mathit{counterexample}.\mathit{hierarchy})=1\).  Assume
the normalized alternative that \(\alpha\) is missing on \(T_{n,q}\)
when \(q>0\), and on \(T_n\) when \(q=0\).  Suppose no
\(y_i\), \(1\leq i\leq\ell-1\), has a color common-missing with
\(y_\ell\) under \(P\)'s coloring.  Assume every proper
exactly-one-uncolored replacement coloring has a missing color at every
vertex.

Suppose the first exact displayed tail edge colored \(\alpha\) is
\(e_j\), where \(2\leq j\leq\ell\).  The edge is supplied as an element
of the original edge type at exact position \(|A|+j-1\), and every
earlier exact tail position is explicitly required not to carry
\(\alpha\).  Then
\noderef{HierarchyEligibleCommonMissingIndices} is nonempty.
The construction yields a replacement presentation for which
\(y_{j-1}\) and \(y_\ell\) have a common missing color.
\end{lemma}
\begin{proof}
Write
\[
 \mu_n=P.\mathit{coloring},\quad
 H=P.\mathit{counterexample}.\mathit{hierarchy},\quad
 A=\noderef{HierarchyStarredStage}(H,q),\quad
 \ell=\noderef{HierarchyFullTailLength}(C),
\]
where \(q\) is the last stored index of \(H\).  Unfold
\noderef{SameGoodHierarchy} in \(P.\mathit{sameGood}\).  Its
\noderef{SameHierarchyStages} clause identifies the original and current
divider counts, base cores, stages, and hence their exact last starred
stage \(A\).  Therefore the original and current full-tail length and
each \noderef{HierarchyFullTailPrefix} agree.  Put
\[
 X=T(y_{j-1})=\noderef{HierarchyFullTailPrefix}(C,j-1),\qquad
 B=T(y_{\ell-1})=\noderef{HierarchyFullTailPrefix}(C,\ell-1),
\]
and let \(x\) and \(y\) be respectively the last displayed vertices of
\(X\) and \(T\).

Apply \noderef{HierarchyCounterexamplePenultimateElementary} to
\(P.\mathit{minimumContext}\).  It gives elementarity of
\(V(B)=\noderef{TreeSequenceVertexSet}(B)\).  The hierarchy inclusions and
the nested and terminal-stage fields of \noderef{TashkinovHierarchy},
together with the definition of \(\ell\), give \(A\preccurlyeq T\).  Since
\(2\le j\le\ell\), arithmetic in the two prefix definitions and the
entrywise clauses of \noderef{TreeSequencePrefix} give
\[
 A\prec X\preccurlyeq B\prec T,\qquad
 |X|=|A|+j-1,\qquad e_j=T.\mathit{edge}(|X|).                \tag{1}
\]
Every relation here is an exact
\noderef{TreeSequenceInitialSegment}, so it retains each individual edge
identity and displayed vertex.

The coloring \(\mu_n\) is proper and leaves exactly \(e_0\) uncolored by
the transported \noderef{HierarchyMinimumCounterexampleContext}.  Apply
the stipulated vertex-missing hypothesis and choose
\[
                         \beta\in\overline\mu_n(x).           \tag{2}
\]
Because \(j-1\in\{1,\ldots,\ell-1\}\), the no-common hypothesis at this
index and \(\alpha\in\overline\mu_n(y)\) imply that \(\beta\) is not
missing at \(y\).  Hence
\[
                              \alpha\ne\beta.                \tag{3}
\]

We discharge the order hypothesis of \noderef{ChangeLemma} exactly.
Let \(i_\alpha,i_\beta\) be the totalized first missing positions of
\(\alpha,\beta\) on \(T\), and let \(v_\alpha,v_\beta\) be their displayed
vertices.  The given \(\alpha\)-missing witness on \(A\) shows
\(i_\alpha.\mathit{val}\le |A|\).  The vertex \(x\) is at position
\(|X|>|A|\) and misses \(\beta\).  If an earlier vertex of \(B\) also
missed \(\beta\), it would be a distinct vertex of the elementary set
\(V(B)\).  The pairwise-disjoint missing-set clause in
\noderef{ColoredVertexSetElementary}, applied to that vertex and \(x\),
would then contradict (2).  Thus
\[
 i_\beta.\mathit{val}=|X|,\qquad v_\beta=x,\qquad
 i_\alpha.\mathit{val}<i_\beta.\mathit{val}.                \tag{4}
\]
In particular \(\beta\) is missing nowhere on \(A\).  Whenever
\(|A|\le r<i_\beta.\mathit{val}=|X|=e_j.\mathit{val}\), the supplied
first-use premise says that
\(\mu_n(T.\mathit{edge}(r))\ne\alpha\).  Equations (3)--(4), the two
missing witnesses on \(B\), this exact interval exclusion, and the stated
positive- or zero-divider normalization satisfy the relevant case of
\noderef{ChangeLemma}.  Consequently
\[
 \noderef{TwoColorComponent}(G,\mu_n,\alpha,\beta,
        \operatorname{inl}v_\alpha)
 =
 \noderef{TwoColorComponent}(G,\mu_n,\alpha,\beta,
        \operatorname{inl}x).                               \tag{5}
\]

Let \(Q\) be the full \((\alpha,\beta)\)-component through
\(\operatorname{inl}y\).  We prove that it is disjoint from (5), without
claiming that it avoids all of \(B\).  The vertices
\(v_\alpha,x,y\) are distinct: the first two have the strict position
order (4), both lie in \(B\), and \(y\notin V(B)\) by (1), vertex
injectivity in \noderef{TreeSequence}, and the definition of
\noderef{TreeSequenceVertexSet}.  The pairwise-disjoint missing-set clause
of \noderef{ColoredVertexSetElementary}, applied to the distinct vertices
\(v_\alpha,x\in V(B)\), gives
\(\beta\notin\overline\mu_n(v_\alpha)\) and
\(\alpha\notin\overline\mu_n(x)\): the missing witness (2) supports the
first conclusion, and
\(\alpha\in\overline\mu_n(v_\alpha)\) supports the second.  By
\noderef{MissingColorIncidenceEndpoint}, the incidence-graph degree at
\(v_\alpha\) and \(y\) is at most one because \(\alpha\) is missing
there; applying the same result with the colors reversed gives the same
bound at \(x\), where \(\beta\) is missing.  If \(Q\) equalled (5), it
would contain three distinct degree-at-most-one nodes.  But
\noderef{TwoColorComponentSpanningPath} spans the component by a simple
path and \noderef{TwoColorIncidenceDegreeAtMostTwo} bounds its full
degrees by two, so only its two ends can have degree at most one.  This
is impossible.  The equal-or-disjoint assertion
\noderef{TwoColorComponentsDisjoint} therefore gives
\[
 Q\cap\noderef{TwoColorComponent}
   (G,\mu_n,\alpha,\beta,\operatorname{inl}v_\alpha)
   =\varnothing.                                             \tag{6}
\]
The missing endpoint \(y\), properness, and
\noderef{MissingColorIncidenceEndpoint} also show that \(Q\) is a
possibly singleton \noderef{TwoColorPathComponent}.

Define
\[
 \sigma_n=\noderef{TwoColorComponentSwap}
       (G,\mu_n,\alpha,\beta,\operatorname{inl}y).            \tag{7}
\]
Invoke \noderef{StableChangeCasesOneTwo} now, before treating the possible
stored seed.  Supply elementarity of \(B\), (3), the two old missing
colors on \(B\), the path certificate, (6), the normalization, and every
ambient field of \(P.\mathit{minimumContext}\).  We obtain
\[
              \noderef{HierarchyWeaklyStable}(H,q,\sigma_n) \tag{8}
\]
and the accompanying mod-coloring and prefix-hierarchy conclusions.
Also \noderef{TwoColorComponentSwapProper} makes \(\sigma_n\) proper, and
the swap fixes the value \(\mathsf{none}\), so it still leaves exactly
\(e_0\) uncolored.

We prove only the edge-value equality needed up to the split:
\[
 \sigma_n(T.\mathit{edge}(r))=\mu_n(T.\mathit{edge}(r))
       \quad\text{for }|A|\le r<|X|.                         \tag{9}
\]
The old value is not \(\alpha\), by first use.  If it is \(\beta\), apply
\noderef{HierarchyTailMissingColorExclusion} to the elementary prefix
\(B\), the exact prefix \(X=T(y_{j-1})\), and the missing vertex \(x\).
The lemma says that this can only be the exact zero-divider SE/RE seed
\(f_n\), with \(A=T_n\), \(r=|A|\), and
\(\mu_n(f_n)=\beta\in D_n\).  Its certified
\noderef{TashkinovSeriesStep} is an actual SE or RE constructor.  In the
SE case, \noderef{MaximumDefectiveChoice} makes this exact edge incident
with its selected vertex of \(T_n\); in the RE case,
\noderef{RevisitingExtensionCertificate} makes it incident with its
inside endpoint in \(T_n\).  Thus \(f_n\) is a
\noderef{TreeIncidentEdge} of \(T_n\).  Under this zero-divider non-PE
alternative, (8) unfolds to
\noderef{StableColoring}(T_n,D_n,\mu_n,\sigma_n)\), whose protected
incident-edge clause fixes \(f_n\).  Every old value outside
\(\{\alpha,\beta\}\) is fixed directly by (7).  This proves (9), with no
claim about \(e_j\) or any later edge.

We next establish the vertexwise equality required on \(B\):
\[
 \overline\sigma_n(v)=\overline\mu_n(v)\quad(v\in V(B)).      \tag{10}
\]
If \(\operatorname{inl}v\notin Q\), no changed edge-node can be incident
with it, since adjacency in \noderef{TwoColorIncidenceGraph} would put the
vertex-node in \(Q\).  If \(\operatorname{inl}v\in Q\), disjointness (6)
shows \(v\ne v_\alpha,x\).  Apply the pairwise-disjoint missing-set clause
of \noderef{ColoredVertexSetElementary} first to \(v,v_\alpha\), using
the \(\alpha\)-missing witness at \(v_\alpha\), and then to \(v,x\),
using (2).  It implies that neither \(\alpha\) nor \(\beta\) is missing
at \(v\).  Expanding
\noderef{MissingColors} and \noderef{PresentColors} supplies incident
edges of both colors, and \noderef{IsProperEdgeColoring} makes each
identity unique.  Both edge-nodes lie in \(Q\), so (7) transposes the two
values and leaves both colors present.  Every other color is unchanged.
This proves (10), including the case where \(Q\) meets \(B\) in internal
degree-two vertices.

The same-good stage equalities identify the actual terminal step in the
\noderef{TashkinovSeries}.  Its PE or SE/RE
\noderef{TashkinovSeriesStep} constructor supplies the final closure.
The exact ordered-prefix and generating-coloring fields in
\noderef{ExtendedTashkinovTree} identify \(T\) with an exact ordered
prefix of that step's final series tree and identify the step coloring with
\(\mu_n\).  Unpacking the augmentation-sequence component of the step's
\noderef{TashkinovClosure} gives a route to the final series tree: from
\(T_n\) in the PE case and from the exact adjoined seed \(T_n+f_n\) in the
SE/RE case.  By
\noderef{TashkinovAugmentationSequence}, this route is the finite
reflexive-transitive chain of one-step augmentations under \(\mu_n\).
First restrict that chain to its exact prefix \(T\), and then to the exact
prefix \(X\), using
\noderef{TashkinovAugmentationSequenceRestrict}.  Weak stability (8)
transports the route through \(A\), as follows.  If \(q>0\), its
\noderef{TreeColoringInvariant} clause on \(A\) fixes every displayed
edge of the old prefix route and preserves each missing witness there.

Suppose \(q=0\) and the final kind is PE.  In this case (8) is the
\noderef{CombinedStableColoring} relation on the terminal ladder tree
\(T_n\) and the PE closure \(R_n\).  Apply
\noderef{CombinedStableClosurePreserved}, using the positivity and
terminal-color identity of \(P\)'s actual replacement ETT, the PE
identity, and the properness and exactly-one-uncoloredness already
established for both \(\mu_n\) and \(\sigma_n\).  Its three conclusions
say, respectively, that
\[
 \sigma_n(f)=\mu_n(f)\quad(f\in E(T_n)),\qquad
 \sigma_n(f)=\mu_n(f)\quad(f\in E(R_n)),                    \tag{11a}
\]
and that the same ordered tree \(R_n\) is the exact
\(\sigma_n\)-Tashkinov closure of the recorded supporting prefix of
\(T_n\).  In particular, this last conclusion preserves the actual PE
closure certificate; it is not being used as an unstated assertion that
an arbitrary route through the union is fixed.

The \noderef{HierarchyBase}.parallel field writes
\(A=T_n\oplus R_n\) as an exact \noderef{OrderedTreeUnion}.  Its
left-prefix clause puts every displayed edge of \(A\) before
\(|T_n|\) at the identical position of \(T_n\), and its
new-entry-from-right clause identifies every later displayed edge of
\(A\) with an exact displayed edge identity of \(R_n\).  Thus (11a)
fixes every displayed edge identity of \(A\).  Its vertex-set equation
also gives
\[
                         V(A)=V(T_n)\cup V(R_n).              \tag{11b}
\]
The two vertexwise clauses of \noderef{StableColoring} in the combined
relation (8), one on \(T_n\) and one on \(R_n\), together with (11b),
therefore give
\(\overline\sigma_n(v)=\overline\mu_n(v)\) for every \(v\in V(A)\).

Now take the old restricted augmentation chain from \(T_n\) to \(A\).
At each old \noderef{TashkinovAugmentation}, its exact new edge is a
displayed edge of the terminal tree \(A\), so its value is unchanged by
(11a) and the two ordered-union edge clauses.  Its missing-color witness
is a vertex of the predecessor and hence, by the successive ordered
initial-segment inclusions, a vertex of \(A\); its missing set is
unchanged by the two stable-coloring vertexwise clauses and (11b).
The length and exact initial-segment clauses are independent of the
coloring.  Induction over the finite chain therefore gives an explicit
\(\sigma_n\)-augmentation route from \(T_n\) to \(A\).

If \(q=0\) and the kind is SE or RE, the nonparallel base clause gives
\(A=T_n\), so there is no route before \(A\) to transport.  In all three
cases we have now obtained the required route to \(A\).  For each
subsequent old \noderef{TashkinovAugmentation} from \(A\) toward \(X\),
(9) fixes its individual new-edge value, while its old missing witness
lies in the predecessor, hence in \(X\preccurlyeq B\), where (10)
preserves that missingness.  A second induction over these unit-length
steps gives exactly
\[
\begin{cases}
T_n\longrightarrow X,&\text{in PE},\\
T_n+f_n\longrightarrow X,&\text{in SE/RE},
\end{cases}                                                   \tag{11}
\]
under \(\sigma_n\), retaining the old identity \(f_n\), its singleton
connecting set, and its exact \noderef{TreeSequenceAdjoinEdge} seed.

The same old finite reflexive-transitive chain from
\noderef{TashkinovAugmentationSequence} passes through the exact
intermediate prefix \(X\).  Indeed, each
\noderef{TashkinovAugmentation} increases length by one, so induction on
the chain reaches the unique stage of length \(|X|\); the exact-prefix
field inherited from \noderef{ExtendedTashkinovTree} identifies it with
\(X\).  Splitting the reflexive-transitive chain at that unique stage gives
the original-color suffix
\[
                         X\longrightarrow T.                 \tag{12}
\]
Equation (10) preserves the old \(\alpha\)-missing witness in
\(A\preccurlyeq X\) and the old \(\beta\)-missing witness \(x\).
By \noderef{TreeMissingColors}, each preserved vertex witness is exactly a
witness for membership in the tree-missing-color set.  Hence both
\(\alpha,\beta\) are missing on \(X\) under \(\sigma_n\).              \(\tag{13}\)

We check the stored first-position and Gamma-cutoff conditions separately.
Fix \(\eta\in\noderef{HierarchyStoredColorsAt}(H,q)\).  This stored-set
definition excludes every color missing on \(A\); since \(\alpha\) is
missing there, \(\eta\ne\alpha\).  If \(\eta\ne\beta\), the component
swap changes no edge to or from color \(\eta\), so its missing vertices
and totalized \noderef{HierarchyFirstMissingPosition} are unchanged.  If
\(\eta=\beta\), equation (4) says that \(x\) is its first old missing
vertex, and (10) preserves missingness at \(x\) and nonmissingness at all
earlier vertices.  Thus in every case
\[
 \noderef{HierarchyFirstMissingPosition}(\sigma_n,T,\eta)
 =\noderef{HierarchyFirstMissingPosition}(\mu_n,T,\eta).     \tag{14}
\]

The color \(\beta\) belongs to no last Gamma pair.  Indeed,
\noderef{GoodHierarchy}.avoidance places every last Gamma color in
\noderef{HierarchyStageMissingColors}(H,q).  That stage is \(T_n\) when
\(q=0\) and \(T_{n,q}=A\) when \(q>0\); in both cases its vertex set is
contained in \(A\preccurlyeq B\).  Since \(\beta\) is missing at \(x\),
the pairwise-disjoint missing-set clause of
\noderef{ColoredVertexSetElementary} says that \(\beta\) is missing at no
distinct vertex of \(B\).  Every vertex of the stage occurs in \(A\),
whereas \(x\) occurs at position \(|X|>|A|\), so vertex injectivity in
\noderef{TreeSequence} makes each such vertex distinct from \(x\).
Consequently \(\beta\) is missing at no vertex of the stage and is outside
the entire last Gamma union.

Suppose now that \(\alpha\in C.\mathit{good}.\Gamma^q_\eta\) for some
stored \(\eta\).  The exact first-use edge is \(e_j=T.\mathit{edge}(|X|)\)
and has old color \(\alpha\).  If the old first missing position of
\(\eta\) were greater than \(|X|\), then
\[
 |A|\le |X|<
 \min\{\noderef{HierarchyFirstMissingPosition}(\mu_n,T,\eta),|T|\},
\]
where \(|X|<|T|\) follows from \(X\preccurlyeq B\prec T\).
Unfolding \noderef{HierarchySegmentColors} would put \(\alpha\) in the
last \((q,\eta)\)-segment, contradicting
\noderef{GoodHierarchy}.avoidance for the assumed Gamma membership.
Therefore
\[
 \alpha\in C.\mathit{good}.\Gamma^q_\eta
 \quad\Longrightarrow\quad
 \noderef{HierarchyFirstMissingPosition}(\mu_n,T,\eta).\mathit{val}
 \le |X|.                                                     \tag{15}
\]
Together with the exclusion of \(\beta\), this is exactly the affected
Gamma cutoff.

Set \(\mathcal C_0=P.\mathit{corresponding}\), and write
\(\mathcal E_0\) for the original fixed ETT and
\(\mathcal E_m=\mathcal C_0.\mathit{replacementETT}\) for the
intermediate ETT.  Apply
\noderef{StableChangeCasesOneTwoSplitTail} with replacement input
\(\mathcal C_0.\mathit{replacementInput}\), ETT \(\mathcal E_m\), base
\(P.\mathit{replacementBase}\), counterexample
\(P.\mathit{counterexample}\), and context
\(P.\mathit{minimumContext}\).  Its remaining hypotheses are exactly (1),
(3), the old missing colors on \(B\), the path and component data
(6)--(7), the normalized Cases 1/2 alternative, the fixed exact prefix
(9), the verified route (11), the old suffix (12), split missingness (13),
vertexwise equality (10), stored first-position equality (14), and cutoff
(15).  Retain the returned maximum property, base, hierarchy, good
hierarchy, and same-good data, and denote its returned correspondence
certificate by
\[
 \mathcal C_1:\quad
 \noderef{CorrespondingExtendedTashkinovTree}
   (\mathcal E_m,\sigma_n,T).
\]
Thus \(\mathcal C_1\) is relative to the intermediate ETT
\(\mathcal E_m\), not to \(\mathcal E_0\).

The swap makes \(\beta\) missing at \(y\): originally \(\alpha\) is
missing and \(\beta\) is present there, so properness gives the unique
incident \(\beta\)-edge in \(Q\), which is changed to \(\alpha\).
Equation (10) keeps \(\beta\) missing at \(x\).  Thus the replacement full
tree is non-elementary.  Form the resulting
\noderef{HierarchyCounterexample} from the returned maximum property,
base, hierarchy, and good hierarchy.

We next compose correspondence back to the original ETT.  Construct
\[
 \mathcal C_{10}:\quad
 \noderef{CorrespondingExtendedTashkinovTree}
   (\mathcal E_0,\sigma_n,T)
\]
field by field as follows.  Take its replacement input, final series, and
final ETT from \(\mathcal C_1\), transporting the series' round parameter
along \(\mathcal C_0.\mathit{rungAgreement}\).  Its rung equality is the
transitive composite
\[
 \mathcal C_1.\mathit{rungAgreement}\mathbin{;}
 \mathcal C_0.\mathit{rungAgreement}.
\]
The heterogeneous series-presentation equality of \(\mathcal C_1\)
remains the required one after this cast.

Here are the component equalities explicitly.  For a positive ladder
index of \(\mathcal E_0\), cast it first to the corresponding index of
\(\mathcal E_m\) using the inverse of
\(\mathcal C_0.\mathit{rungAgreement}\).  The new ladder tree equals the
ladder tree in \(\mathcal E_m.\mathit{series}\) by
\(\mathcal C_1.\mathit{ladderTreesAgree}\); the
\(\mathcal C_0.\mathit{seriesPresentation}\) field identifies that series,
with the same dependent index cast, with
\(\mathcal C_0.\mathit{replacementSeries}\); and
\(\mathcal C_0.\mathit{ladderTreesAgree}\) identifies the latter tree with
the ladder tree of \(\mathcal E_0\).  This three-step chain is the
ladder-tree field of \(\mathcal C_{10}\).  At an index in the full
round-plus-one range, the identical chain, using first
\(\mathcal C_1\)'s equality, then
\(\mathcal C_0.\mathit{seriesPresentation}\), and finally the corresponding
equality of \(\mathcal C_0\), gives respectively the stored-set,
connecting-edge-set, and extension-kind fields.  All four chains use the
dependent casts induced by
\(\mathcal C_1.\mathit{rungAgreement}\) followed by
\(\mathcal C_0.\mathit{rungAgreement}\); hence they compare the final
series with the original series, rather than stopping at
\(\mathcal E_m\).  Take final-generating-coloring, replacement properness,
and replacement one-uncoloredness from \(\mathcal C_1\), after the same
round-index casts.

It remains to construct the replacement-stability field of
\(\mathcal C_{10}\).  Fix a positive ladder index \(i\) of
\(\mathcal E_0\), and use the rung cast above for every intermediate
index.  The \(i\)-th relation in
\(\mathcal C_0.\mathit{replacementStable}\) is stability from the original
series coloring to the intermediate series coloring.  Transport
\(\mathcal C_1.\mathit{replacementStable}\) through
\(\mathcal C_0.\mathit{seriesPresentation}\) and the ladder-tree equality;
it is stability from that same intermediate coloring to the new series
coloring on the original ladder tree.  The stable vertexwise equality in
the first relation, together with the raw stored-set and ladder-tree
equalities just proved, gives equality of the two computed
\noderef{TashkinovStoredColors} sets, so both relations use the same tree
and protected set.

The subtype hypotheses required for transitivity in \noderef{sc2} are
available at this exact index.  The intermediate coloring is proper and
leaves exactly \(e_0\) uncolored by the two corresponding fields of
\(\mathcal C_0\), and the new coloring has those properties by the two
fields of \(\mathcal C_1\).  For the original coloring, inspect the
certified positive \noderef{TashkinovSeriesStep}; its initial constructor
cannot occur at \(i+1\).  In the SE constructor, the current-coloring
equality identifies the coloring with the selected coloring, whose
properness and one-uncoloredness are fields of
\noderef{MaximumDefectiveChoice}.  In the PE constructor they are the
current-proper and current-one-uncolored fields.  In the RE constructor
the current-coloring equality identifies the coloring with the preceding
one, for which the constructor supplies previous-proper and
previous-one-uncolored.  Put these three colorings into the subtype in
\noderef{sc2}; its transitivity clause composes the original-to-intermediate
and intermediate-to-new stable relations.  This is exactly the
\(i\)-th replacement-stability field of \(\mathcal C_{10}\), completing
the composite certificate.

The returned maximum property, base, hierarchy, good hierarchy, and the
new counterexample remain indexed by
\(\mathcal C_1.\mathit{replacementETT}\), which is definitionally the final
ETT in \(\mathcal C_{10}\); likewise both certificates use
\(\mathcal C_1.\mathit{replacementInput}\) after the cast above.  Unpack
the fields of \(\mathcal C_{10}\).  Its composite rung agreement
transports the lower-rung package.  Its composite replacement-stability
at the terminal ladder tree transports the \(2n+11\) estimate, and its
properness and one-uncoloredness fields supply those coloring hypotheses.
The returned same-good relation identifies the intermediate and final
divider counts and stages.  Together with \(P.\mathit{sameGood}\), it
therefore transports the fewer-divider field and
\noderef{HierarchyCounterexampleMeasure}, so the original rung-local
minimality field applies.  These data give every field of the replacement
\noderef{HierarchyMinimumCounterexampleContext}, now with precisely
\(\mathcal C_{10}.\mathit{replacementInput}\), its final ETT, base, and
counterexample as indices.

Finally compose the two \noderef{SameGoodHierarchy} witnesses field by
field.  In their \noderef{SameHierarchyStages} components, compose the
core equalities, divider-count equalities, heterogeneous stage and
divider-position equalities, PE equivalences, and, in the PE case, closure
equalities.  At each divider index cast along the composite count equality,
compose the two stored-set equalities and then the two Gamma-set
equalities for every color in that common stored set.  This yields
same-good data from the original \(C.\mathit{good}\) directly to the final
good hierarchy.  Using \(\mathcal C_{10}\), rather than \(\mathcal C_1\),
as the corresponding field, all the preceding data construct a genuine
replacement presentation \(P'\) relative to the original fixed ETT.

Finally \(1\le j-1\le\ell-1\), and under \(P'\) the color \(\beta\) is
missing at both \(x=y_{j-1}\) and \(y=y_\ell\).  The witnesses
\(j-1,P',\beta\) therefore satisfy the defining filter of
\noderef{HierarchyEligibleCommonMissingIndices}, proving the claim.
\end{proof}
